\label{gu:app:functional}
\subsection{Tensor evaluation and closed-subgroup duality}
\begin{lemma}[Tensor evaluation]\label{gu:app:tensor}\label{fun:tensor}
Let $A$ and $B$ be locally compact totally disconnected groups,
$\mathscr E$ a smooth $A\times B$-module, $\pi$ an admissible smooth
$A$-module, and $V$ a smooth $B$-module. Contraction gives a natural
isomorphism
\begin{equation}\label{gu:eq:tensor-lemma}
 \Hom_B((\mathscr E\otimes\pi^\vee)_A,V)
 \simeq\Hom_{A\times B}(\mathscr E,\pi\boxtimes V).
\end{equation}
All tensor products here are algebraic.
\end{lemma}
\begin{proof}
Let $b:\mathscr E\times\pi^\vee\to V$ represent a map on the left.
For $e\in\mathscr E$, choose a compact open subgroup $J\subset A$
fixing $e$. Diagonal invariance makes $b(e,-)$ factor through
$(\pi^\vee)_J=(\pi^\vee)^J$, where the equality is given by compact
averaging. Its pairing with $\pi^J$ is perfect: any functional on
$\pi^J$ extends to a smooth functional on $\pi$ by composition with
the averaging projection. Admissibility makes both spaces finite
dimensional. For dual bases $e_i\in\pi^J$ and
$e_i^\vee\in(\pi^\vee)^J$, the corresponding tensor is
\[
 F(e)=\sum_i e_i\otimes b(e,e_i^\vee)\in\pi\otimes V.
\]
It is characterized by $(\ell\otimes1)F(e)=b(e,\ell)$ for all
$\ell\in\pi^\vee$. This proves independence of $J$ and the bases,
linearity, and equivariance. In particular,
$b(ae,\ell)=b(e,\pi^\vee(a^{-1})\ell)$ gives the $\pi(a)$
action. Contraction is the inverse. 
\end{proof}

For irreducible \(\pi\), this identifies the coinvariant construction
with the usual big theta lift. To see the maximal-isotypic
statement explicitly, put \(T=(\mathscr E\otimes\pi^\vee)_A\).
The identity of \(T\) gives a map \(q:\mathscr E\to\pi\otimes T\).
The target is semisimple and isotypic, so its submodules have the
form \(\pi\otimes T_0\). If this is the image of \(q\), all
contractions \([e\otimes\ell]\) belong to \(T_0\). Those classes
span \(T\), proving surjectivity. Lemma~\ref{gu:app:tensor} supplies
the universal property.

\begin{lemma}[Closed-subgroup duality]\label{gu:app:duality}
Let $D$ be a closed unimodular subgroup of a unimodular locally compact
totally disconnected group $L$, and let $V$ be an admissible smooth
$L$-module. A nonzero invariant measure on $D\backslash L$ gives
\begin{equation}\label{gu:eq:closed-duality}
 \Hom_L(\cInd_D^L\mathbf1,V)
 \simeq\Hom_D(V^\vee,\mathbf1).
\end{equation}
\end{lemma}
\begin{proof}
Integration against the quotient measure identifies
\begin{equation}\label{gu:eq:compact-dual}
 (\cInd_D^L\mathbf1)^\vee=\Ind_D^L\mathbf1.
\end{equation}
Indeed, for a compact open $J\subset L$, the $J$-orbits on
$D\backslash L$ are compact open. A $J$-invariant distribution is a
scalar multiple of invariant measure on each orbit. Smooth
distributions are therefore right-smooth locally constant densities,
with no compact-support condition; every such density pairs with
compactly supported functions.

Smooth reflexivity of $V$ and ordinary smooth-induction reciprocity give
\[
 \begin{aligned}
 \Hom_L(\cInd_D^L\mathbf1,V)
 &\simeq\Hom_L(V^\vee,\Ind_D^L\mathbf1)\\
 &\simeq\Hom_D(V^\vee,\mathbf1).
 \end{aligned}
\]
For inverse transposition, evaluating on a compactly induced vector
fixed by a compact open subgroup gives a smooth functional on
$V^\vee$, hence a vector of $V$. The last map is evaluation at $1$; its inverse sends a
$D$-invariant functional $\ell$ to
$v^\vee\mapsto[k\mapsto\ell(kv^\vee)]$.
\end{proof}

\subsection{Fiber relations and integration}
\begin{lemma}\label{gu:app:fibers}
Let $N$ act smoothly on a vector space $E_0$ by $T$, and let
$\alpha_c$ be automorphisms of $N$, for $c\in F^\times$, such that
$c\mapsto T_{\alpha_c(n)}v$ is locally constant for fixed $n,v$.
Give $C_c^\infty(F^\times,E_0)$ the action
$[nF](c)=T_{\alpha_c(n)}F(c)$. Then pointwise projection induces
\[
 C_c^\infty(F^\times,E_0)_N
 \simeq C_c^\infty(F^\times,(E_0)_N).
\]
If the remaining multiplier action is pure translation
$F(c)\mapsto F(ct)$ or $F(c/t)$, its coinvariant quotient is
$(E_0)_N$, with quotient map integration against $d^\times c$.
\end{lemma}
\begin{proof}
Surjectivity follows by lifting the finitely many values of a locally
constant compactly supported function. For the kernel, partition the
support into finitely many compact open pieces and express each value
as a finite sum $(T_w-1)v$. At a point $c_0$ put
$n=\alpha_{c_0}^{-1}(w)$. Local constancy gives
$T_{\alpha_c(n)}v=T_wv$ on a neighborhood of $c_0$. Multiplication by
its characteristic function realizes this as a global $n$-relation.
A finite cover and disjoint refinement give the whole kernel; the
reverse inclusion is pointwise.

For translation, write $F=\sum_i1_{c_iJ}m_i$ for a common compact open
subgroup $J\subset F^\times$. Its integral is
$\operatorname{vol}(J)\sum_i m_i$. If this is zero, subtracting one
fixed coset expresses $F$ as a sum of translation differences.
Surjectivity follows from
$m\mapsto\operatorname{vol}(J)^{-1}1_Jm$. Both orientations have the
same relation span. All integrals are finite vector-valued sums.
\end{proof}

\subsection{Matching and product coinvariants}
\begin{lemma}\label{gu:app:comparison}\label{fun:reverse-comparison}\label{fun:forward-comparison}
Let $\varpi$ be a smooth representation of the matching group $R$ on
$\mathscr S$, and let $\Omega=\cInd_R^{G_\theta\times G'_\theta}\varpi$, realized
by section values $f(c)=f(q_c)$. For irreducible admissible $\tau$ of
$G_\theta$ and $\Pi$ of $G'_\theta$, put
\[
 V_\tau=(\mathscr S\otimes\tau^\vee)_{G_\theta^1},\qquad
 M_\Pi=(\mathscr S\otimes\Pi^\vee)_{(G'_\theta)^1}.
\]
Their matching actions are
\begin{align}
 h[v\otimes\ell_\tau]
  &=[\varpi(r_t,h)v\otimes\tau^\vee(r_t)\ell_\tau],
     &&t=\lambda(h),\label{gu:eq:forward-matching-action}\\
 g[v\otimes\ell]
  &=[\varpi(g,d_s)v\otimes\Pi^\vee(d_s)\ell],
     &&s=\lambda(g).\label{gu:eq:reverse-matching-action}
\end{align}
Integration gives $G'_\theta$- and $G_\theta$-equivariant isomorphisms:
\begin{align}
 (\Omega\otimes\tau^\vee)_{G_\theta}&\simeq V_\tau,
 &f\otimes\ell_\tau&\longmapsto
 \int[\varpi(e_c)^{-1}f(c)\otimes\tau^\vee(r_c^{-1})\ell_\tau]\,d^\times c,
 \label{gu:eq:forward-comparison}\\
 (\Omega\otimes\Pi^\vee)_{G'_\theta}&\simeq M_\Pi,
 &f\otimes\ell&\longmapsto
 \int[f(c)\otimes\Pi^\vee(d_c)\ell]\,d^\times c.
 \label{gu:asm:quotient-integral}
\end{align}
For $\varpi=\omega_{R,0}$, the first integrand uses $D_c^{-1}$.
\end{lemma}
\begin{proof}
The matching actions are independent of the chosen matching elements:
two choices differ by the group whose coinvariants are being taken.
Normality preserves the relations, and the multiplicative sections give
the group law and smoothness.

For the forward map, $\gamma(g)=\varpi(g,d_{\lambda(g)})$ is a
representation of $G_\theta$, and its induced action is
$f(c)\mapsto\gamma(g)f(c/\lambda(g))$. Lemma~\ref{gu:app:fibers}, initially
with constant automorphism, gives $G_\theta^1$-coinvariants
$C_c^\infty(F^\times,V_\tau)$. The residual action has fiber operator
$\varpi(e_s)\otimes\tau^\vee(r_s)$. Applying its inverse at $c$
turns this into pure translation. The same lemma proves that the
forward integral has exactly the $G_\theta$-relation space as kernel and is
onto. For $G'_\theta$-equivariance, substitute $q=ct$ and use
\[
 \varpi(e_{q/t})^{-1}\varpi(1,d_{q/t}hd_q^{-1})
 =\varpi(r_t,h)\varpi(e_q)^{-1},\qquad
 r_{q/t}^{-1}=r_t r_q^{-1}.
\]
This transports the action to \eqref{gu:eq:forward-matching-action}.

For the reverse map, identify the tensor space with
$C_c^\infty(F^\times,E_0)$, where $E_0=\mathscr S\otimes\Pi^\vee$,
and change fibers by
\begin{equation}\label{gu:asm:fiber-change}
 (\mathcal BF)(c)=(1\otimes\Pi^\vee(d_c))F(c).
\end{equation}
This and its inverse preserve local constancy and compact support.
For $u\in (G'_\theta)^1$ the new fiber action is
$T_{d_cud_c^{-1}}$, where $T_u=\varpi(1,u)\otimes\Pi^\vee(u)$.
Lemma~\ref{gu:app:fibers} applies with $\alpha_c=\operatorname{Ad}(d_c)$.
After taking these relations, $d_t$ acts by $F(c)\mapsto F(ct)$.
The lemma therefore proves that the reverse integral is onto with
the $G'_\theta$-relation space as kernel.

The equivariance directions can also be read directly. Put
$u(c,h)=d_chd_{ct}^{-1}\in (G'_\theta)^1$. Diagonal $h$-translation changes the
reverse integrand to
\[
 [\varpi(1,u(c,h))f(ct)\otimes\Pi^\vee(d_ch)\ell]
 =[f(ct)\otimes\Pi^\vee(d_{ct})\ell],
\]
using the coinvariant relation. Multiplicative Haar invariance proves
$G'_\theta$-invariance. For $g\in G_\theta$, the identity
$q_c(g,1)=(g,d_s)q_{c/s}$ and substitution $c=ys$ give
\[
 \int[\varpi(g,d_s)f(y)\otimes
             \Pi^\vee(d_s)\Pi^\vee(d_y)\ell]\,d^\times y,
\]
which is the action in \eqref{gu:eq:reverse-matching-action}.
Unimodularity of $R$ and $G_\theta\times G'_\theta$ makes the quotient
measure multiplicative Haar measure.
\end{proof}
